%% Generated by tools/gen_error_codes.py from docs/errors/ of the compiler;
%% edit the compiler's pages or tools/error_themes.txt, then rerun it.

\clearpage
\section{Unions}
\label{sec:codes:unions}

\begin{longtable}{@{}l p{0.78\linewidth} r@{}}
\toprule Code & Message & Page\\ \midrule \endhead
\hyperref[sec:codes:e4028]{E4028} & type \errhole{} cannot be embedded within a union type & \pageref{sec:codes:e4028}\\
\hyperref[sec:codes:e4306]{E4306} & over is prohibited on a union of records & \pageref{sec:codes:e4306}\\
\hyperref[sec:codes:e4307]{E4307} & union \errhole{} declares the member \errhole{} twice & \pageref{sec:codes:e4307}\\
\hyperref[sec:codes:e4308]{E4308} & union \errhole{} declares no member & \pageref{sec:codes:e4308}\\
\hyperref[sec:codes:e4309]{E4309} & a union cannot be a member of another union & \pageref{sec:codes:e4309}\\
\hyperref[sec:codes:e4310]{E4310} & the member of a union must be a record or a class type, not \errhole{} & \pageref{sec:codes:e4310}\\
\hyperref[sec:codes:e4311]{E4311} & union \errhole{} mixes record and class members & \pageref{sec:codes:e4311}\\
\hyperref[sec:codes:e4312]{E4312} & union \errhole{} has no method named \errhole{} & \pageref{sec:codes:e4312}\\
\hyperref[sec:codes:e4313]{E4313} & union \errhole{} cannot be converted to \errhole{} & \pageref{sec:codes:e4313}\\
\hyperref[sec:codes:e4314]{E4314} & union \errhole{} does not list the member \errhole{} & \pageref{sec:codes:e4314}\\
\hyperref[sec:codes:e4315]{E4315} & \errhole{} is not a union type & \pageref{sec:codes:e4315}\\
\hyperref[sec:codes:e4317]{E4317} & union \errhole{} does not implement trait \errhole{} & \pageref{sec:codes:e4317}\\
\bottomrule
\end{longtable}

\errorcode{E4028}{type \errhole{} cannot be embedded within a union type}
\label{sec:codes:e4028}

Raised during \textbf{validation}, from \texttt{ValidateErrorMessage::\allowbreak{}CONTAIN\_\allowbreak{}ALIASABLE\_\allowbreak{}TYPE} (\textit{src/\allowbreak{}ymirc/\allowbreak{}semantic/\allowbreak{}validator/\allowbreak{}errors.yr}).

\subsubsection*{What it means}

A type that can be aliased -- one that borrows memory -- was placed in a union. The fields of a union share their storage, so the compiler cannot know which of them owns the memory at a given moment.

\subsubsection*{Example}

\textit{test\_\allowbreak{}resources/\allowbreak{}union/\allowbreak{}test5.yr}:

%% check: skip
\begin{lstlisting}[style=coloredverbatimError]
@overlaid
record A {
    let a : [i32];
}
\end{lstlisting}

\begin{lstlisting}[style=bashVerb, breaklines=true, escapechar=~]
Error[E4028] : type [i32] cannot be embedded within a union type
 --> test_resources/union/test5.yr:(3,5)
 3  ~\lstbox{\textSFxi{}}~     let a : [i32];
    ~\lstbox{\textSFv{}}~     ^^^     ^
    ~\lstbox{\textSFxi{}}~ Note : union cannot contain aliasable types, with the exceptions of pointer types
    ~\lstbox{\textSFxi{}}~  --> test_resources/union/test5.yr:(2,8)
    ~\lstbox{\textSFxi{}}~  2  ~\lstbox{\textSFxi{}}~ record A {
    ~\lstbox{\textSFxi{}}~     ~\lstbox{\textSFv{}}~        ^
    ~\lstbox{\textSFxi{}}~     ~\lstbox{\textSFii{}}~~\lstbox{\textSFx{}}~~\lstbox{\textSFx{}}~~\lstbox{\textSFx{}}~~\lstbox{\textSFx{}}~~\lstbox{\textSFx{}}~~\lstbox{\textSFx{}}~~\lstbox{\textSFx{}}~
    ~\lstbox{\textSFii{}}~~\lstbox{\textSFx{}}~~\lstbox{\textSFx{}}~~\lstbox{\textSFx{}}~~\lstbox{\textSFx{}}~~\lstbox{\textSFx{}}~~\lstbox{\textSFvii{}}~~\lstbox{\textSFx{}}~
\end{lstlisting}

\subsubsection*{How to fix it}

Store the borrowing type outside the union, and keep only plain values in it.

\errorcode{E4306}{over is prohibited on a union of records}
\label{sec:codes:e4306}

Raised during \textbf{validation}, from \texttt{ValidateErrorMessage::\allowbreak{}UNION\_\allowbreak{}BASE\_\allowbreak{}ON\_\allowbreak{}RECORD} (\textit{src/\allowbreak{}ymirc/\allowbreak{}semantic/\allowbreak{}validator/\allowbreak{}errors.yr}).

\subsubsection*{What it means}

A union of records is perfectly legal -- what is rejected here is the \tokennolst{over} on it. \tokennolst{over B} states that a vtable suitable for every member already exists, so that a method of \tokennolst{B} can be called on the union whichever member it holds. Records have no vtable and no inheritance, so there is nothing for a union of records to declare a base of. The note names the member that makes it one.

A union of records still gets the methods its members have in common: the compiler synthesizes a table of the methods every member declares with an identical signature, without a base being named.

\subsubsection*{Example}

\textit{test\_\allowbreak{}resources/\allowbreak{}type\_\allowbreak{}union/\allowbreak{}test10.yr}:

%% check: skip
\begin{lstlisting}[style=coloredverbatimError]
record A {
    let a : i32;
}

record B {
    let b : i32;
}

class Base {
    pub self () {}
}

union U over Base
| A
| B;
\end{lstlisting}

\begin{lstlisting}[style=bashVerb, breaklines=true, escapechar=~]
Error[E4306] : over is prohibited on a union of records
 --> test_resources/type_union/test10.yr:(15,14)
15  ~\lstbox{\textSFxi{}}~ union U over Base
    ~\lstbox{\textSFv{}}~              ^^^^
16  ~\lstbox{\textSFxi{}}~ | A
    ~\lstbox{\textSFv{}}~   -  note: test10::A is a record type
    ~\lstbox{\textSFxi{}}~
    = when validating test10::U -- (test_resources/type_union/test10.yr:(15,7))
    ~\lstbox{\textSFii{}}~~\lstbox{\textSFx{}}~~\lstbox{\textSFx{}}~~\lstbox{\textSFx{}}~~\lstbox{\textSFx{}}~~\lstbox{\textSFx{}}~~\lstbox{\textSFx{}}~~\lstbox{\textSFx{}}~
\end{lstlisting}

\subsubsection*{How to fix it}

Drop the \tokennolst{over}, or make the members classes descending from the base.

\errorcode{E4307}{union \errhole{} declares the member \errhole{} twice}
\label{sec:codes:e4307}

Raised during \textbf{validation}, from \texttt{ValidateErrorMessage::\allowbreak{}UNION\_\allowbreak{}DUPLICATE\_\allowbreak{}MEMBER} (\textit{src/\allowbreak{}ymirc/\allowbreak{}semantic/\allowbreak{}validator/\allowbreak{}errors.yr}).

\subsubsection*{What it means}

A member is identified at runtime by its position in the member list, so two members naming the same type are two names for one case: nothing can ever produce the second, and the second of two arms matching it is unreachable. The duplicate is reported rather than dropped, the note pointing at the first occurrence.

\subsubsection*{Example}

\textit{test\_\allowbreak{}resources/\allowbreak{}type\_\allowbreak{}union/\allowbreak{}test7.yr}:

%% check: skip
\begin{lstlisting}[style=coloredverbatimError]
union U
| A
| B
| A;
\end{lstlisting}

\begin{lstlisting}[style=bashVerb, breaklines=true, escapechar=~]
Error[E4307] : union test7::U declares the member test7::A twice
 --> test_resources/type_union/test7.yr:(14,3)
12  ~\lstbox{\textSFxi{}}~ | A
    ~\lstbox{\textSFv{}}~   -
13  ~\lstbox{\textSFxi{}}~ | B
14  ~\lstbox{\textSFxi{}}~ | A;
    ~\lstbox{\textSFv{}}~   ^
    ~\lstbox{\textSFxi{}}~
    = when validating test7::U -- (test_resources/type_union/test7.yr:(11,7))
    ~\lstbox{\textSFii{}}~~\lstbox{\textSFx{}}~~\lstbox{\textSFx{}}~~\lstbox{\textSFx{}}~~\lstbox{\textSFx{}}~~\lstbox{\textSFx{}}~~\lstbox{\textSFx{}}~~\lstbox{\textSFx{}}~
\end{lstlisting}

\subsubsection*{How to fix it}

Remove the repeated member.

\errorcode{E4308}{union \errhole{} declares no member}
\label{sec:codes:e4308}

Raised during \textbf{validation}, from \texttt{ValidateErrorMessage::\allowbreak{}UNION\_\allowbreak{}EMPTY} (\textit{src/\allowbreak{}ymirc/\allowbreak{}semantic/\allowbreak{}validator/\allowbreak{}errors.yr}).

\subsubsection*{What it means}

A union holds exactly one of its members. With no member there is nothing to build such a value from, no arm that could match it, and no default value to fall back on, so the type is uninhabited and the declaration is rejected.

A union of a \textit{single} member is accepted: written by hand it is a redundant alias, but a template union can specialize down to one member.

\subsubsection*{Example}

\textit{test\_\allowbreak{}resources/\allowbreak{}type\_\allowbreak{}union/\allowbreak{}test6.yr}:

%% check: skip
\begin{lstlisting}[style=coloredverbatimError]
union U;
\end{lstlisting}

\begin{lstlisting}[style=bashVerb, breaklines=true, escapechar=~]
Error[E4308] : union test6::U declares no member
 --> test_resources/type_union/test6.yr:(3,7)
 3  ~\lstbox{\textSFxi{}}~ union U;
    ~\lstbox{\textSFv{}}~       ^
    ~\lstbox{\textSFxi{}}~
    = when validating test6::U -- (test_resources/type_union/test6.yr:(3,7))
    ~\lstbox{\textSFii{}}~~\lstbox{\textSFx{}}~~\lstbox{\textSFx{}}~~\lstbox{\textSFx{}}~~\lstbox{\textSFx{}}~~\lstbox{\textSFx{}}~~\lstbox{\textSFx{}}~~\lstbox{\textSFx{}}~
\end{lstlisting}

\subsubsection*{How to fix it}

List the members the union can hold.

\errorcode{E4309}{a union cannot be a member of another union}
\label{sec:codes:e4309}

Raised during \textbf{validation}, from \texttt{ValidateErrorMessage::\allowbreak{}UNION\_\allowbreak{}MEMBER\_\allowbreak{}IN\_\allowbreak{}UNION} (\textit{src/\allowbreak{}ymirc/\allowbreak{}semantic/\allowbreak{}validator/\allowbreak{}errors.yr}).

\subsubsection*{What it means}

A union listed as the member of another union is neither flattened nor nested. Flattening would make the member list of the outer union depend on a declaration written elsewhere, and a match over it would have no member to name, the inner union not being one of its cases.

\subsubsection*{Example}

\textit{test\_\allowbreak{}resources/\allowbreak{}type\_\allowbreak{}union/\allowbreak{}test8.yr}:

%% check: skip
\begin{lstlisting}[style=coloredverbatimError]
union U
| A
| B;

union V
| A
| U;
\end{lstlisting}

\begin{lstlisting}[style=bashVerb, breaklines=true, escapechar=~]
Error[E4309] : a union cannot be a member of another union
 --> test_resources/type_union/test8.yr:(17,3)
17  ~\lstbox{\textSFxi{}}~ | U;
    ~\lstbox{\textSFv{}}~   ^
 --> test_resources/type_union/test8.yr:(11,7)
11  ~\lstbox{\textSFxi{}}~ union U
    ~\lstbox{\textSFv{}}~       -
    ~\lstbox{\textSFxi{}}~
    = when validating test8::V -- (test_resources/type_union/test8.yr:(15,7))
    ~\lstbox{\textSFii{}}~~\lstbox{\textSFx{}}~~\lstbox{\textSFx{}}~~\lstbox{\textSFx{}}~~\lstbox{\textSFx{}}~~\lstbox{\textSFx{}}~~\lstbox{\textSFx{}}~~\lstbox{\textSFx{}}~
\end{lstlisting}

\subsubsection*{How to fix it}

List the members of the inner union directly.

\errorcode{E4310}{the member of a union must be a record or a class type, not \errhole{}}
\label{sec:codes:e4310}

Raised during \textbf{validation}, from \texttt{ValidateErrorMessage::\allowbreak{}UNION\_\allowbreak{}MEMBER\_\allowbreak{}NOT\_\allowbreak{}RECORD\_\allowbreak{}NOR\_\allowbreak{}CLASS} (\textit{src/\allowbreak{}ymirc/\allowbreak{}semantic/\allowbreak{}validator/\allowbreak{}errors.yr}).

\subsubsection*{What it means}

The members of a union are the named types it can hold, and only a record or a class is one. A scalar, a slice or an entity is rejected -- an entity carrying a note saying so, since it otherwise reads like a record.

\subsubsection*{Example}

\textit{test\_\allowbreak{}resources/\allowbreak{}type\_\allowbreak{}union/\allowbreak{}test12.yr}:

%% check: skip
\begin{lstlisting}[style=coloredverbatimError]
union U
| i32
| [c8];
\end{lstlisting}

\begin{lstlisting}[style=bashVerb, breaklines=true, escapechar=~]
Error[E4310] : the member of a union must be a record or a class type, not i32
 --> test_resources/type_union/test12.yr:(4,3)
 4  ~\lstbox{\textSFxi{}}~ | i32
    ~\lstbox{\textSFv{}}~   ^^^
    ~\lstbox{\textSFxi{}}~
    = when validating test12::U -- (test_resources/type_union/test12.yr:(3,7))
    ~\lstbox{\textSFii{}}~~\lstbox{\textSFx{}}~~\lstbox{\textSFx{}}~~\lstbox{\textSFx{}}~~\lstbox{\textSFx{}}~~\lstbox{\textSFx{}}~~\lstbox{\textSFx{}}~~\lstbox{\textSFx{}}~
\end{lstlisting}

\subsubsection*{How to fix it}

Wrap the value in a record, and list that record as the member.

\errorcode{E4311}{union \errhole{} mixes record and class members}
\label{sec:codes:e4311}

Raised during \textbf{validation}, from \texttt{ValidateErrorMessage::\allowbreak{}UNION\_\allowbreak{}MIXED\_\allowbreak{}MEMBERS} (\textit{src/\allowbreak{}ymirc/\allowbreak{}semantic/\allowbreak{}validator/\allowbreak{}errors.yr}).

\subsubsection*{What it means}

The members of a union share one overlay, every member sitting at the same offset. A record is held there as an inline value and a class as a pointer, so mixing the two would put the value semantics of one and the reference semantics of the other in a single slot. A union is therefore over records, or over classes, never both. The note points at the member that set the kind.

\subsubsection*{Example}

\textit{test\_\allowbreak{}resources/\allowbreak{}type\_\allowbreak{}union/\allowbreak{}test9.yr}:

%% check: skip
\begin{lstlisting}[style=coloredverbatimError]
record A {
    let a : i32;
}

class B {
    pub self () {}
}

union U
| A
| B;
\end{lstlisting}

\begin{lstlisting}[style=bashVerb, breaklines=true, escapechar=~]
Error[E4311] : union test9::U mixes record and class members
 --> test_resources/type_union/test9.yr:(13,3)
12  ~\lstbox{\textSFxi{}}~ | A
    ~\lstbox{\textSFv{}}~   -
13  ~\lstbox{\textSFxi{}}~ | B;
    ~\lstbox{\textSFv{}}~   ^
    ~\lstbox{\textSFxi{}}~
    = when validating test9::U -- (test_resources/type_union/test9.yr:(11,7))
    ~\lstbox{\textSFii{}}~~\lstbox{\textSFx{}}~~\lstbox{\textSFx{}}~~\lstbox{\textSFx{}}~~\lstbox{\textSFx{}}~~\lstbox{\textSFx{}}~~\lstbox{\textSFx{}}~~\lstbox{\textSFx{}}~
\end{lstlisting}

\subsubsection*{How to fix it}

Make every member a record, or every member a class.

\errorcode{E4312}{union \errhole{} has no method named \errhole{}}
\label{sec:codes:e4312}

Raised during \textbf{validation}, from \texttt{ValidateErrorMessage::\allowbreak{}UNION\_\allowbreak{}NO\_\allowbreak{}METHOD} (\textit{src/\allowbreak{}ymirc/\allowbreak{}semantic/\allowbreak{}validator/\allowbreak{}errors.yr}).

\subsubsection*{What it means}

A union holds one of its members, which one being known only at runtime, so a method call on it has to be dispatched whichever member it holds. What it answers to is therefore not the union of the methods of its members but their intersection: a method every member declares publicly, under the same name and with the same signature, the return type included. A union declaring a base answers to the public methods of that base instead, every member descending from it.

An operator applied to a union is such a call too, and the method it names is the one reported: \tokennolst{u [i]} calls \tokennolst{opIndex}, \tokennolst{u [i] = v} calls \tokennolst{opIndexAssign}, \tokennolst{u == v} calls \tokennolst{opEquals}, and \tokennolst{u \textless{} v} calls \tokennolst{opCmp}.

The notes name, member by member, what each declares under that name: the prototype it declares, that it declares none, or that the one it declares is private. Comparing them shows which member keeps the method out of the union.

\subsubsection*{Example}

\textit{test\_\allowbreak{}resources/\allowbreak{}type\_\allowbreak{}union/\allowbreak{}test23.yr}:

%% check: skip
\begin{lstlisting}[style=coloredverbatimError]
class A {
    pub self () {}
    pub fn foo (self)-> i32 { 1 }
}

class B {
    pub self () {}
}

union U
| A
| B;

fn bar (u : &U)-> i32 {
    u.foo ()
}
\end{lstlisting}

\begin{lstlisting}[style=bashVerb, breaklines=true, escapechar=~]
Error[E4312] : union &(test23::U) has no method named foo
 --> test_resources/type_union/test23.yr:(17,7)
17  ~\lstbox{\textSFxi{}}~     u.foo ()
    ~\lstbox{\textSFv{}}~       ^^^
 --> test_resources/type_union/test23.yr:(5,12)
 5  ~\lstbox{\textSFxi{}}~     pub fn foo (self)-> i32 { 1 }
    ~\lstbox{\textSFv{}}~            ---  note: the member test23::A declares test23::A::foo ()-> i32
   ...
14  ~\lstbox{\textSFxi{}}~ | B;
    ~\lstbox{\textSFv{}}~   -  note: the member test23::B declares no method named foo
    ~\lstbox{\textSFxi{}}~
    = when validating test23::bar -- (test_resources/type_union/test23.yr:(16,4))
    ~\lstbox{\textSFii{}}~~\lstbox{\textSFx{}}~~\lstbox{\textSFx{}}~~\lstbox{\textSFx{}}~~\lstbox{\textSFx{}}~~\lstbox{\textSFx{}}~~\lstbox{\textSFx{}}~~\lstbox{\textSFx{}}~
\end{lstlisting}

\subsubsection*{How to fix it}

Declare the method in every member, publicly and with the same signature, or reach the member through a match and call it there.

\errorcode{E4313}{union \errhole{} cannot be converted to \errhole{}}
\label{sec:codes:e4313}

Raised during \textbf{validation}, from \texttt{ValidateErrorMessage::\allowbreak{}UNION\_\allowbreak{}NARROWING} (\textit{src/\allowbreak{}ymirc/\allowbreak{}semantic/\allowbreak{}validator/\allowbreak{}errors.yr}).

\subsubsection*{What it means}

Conversion into a union is one way. A member widens into the union listing it wherever a conversion is allowed -- a variable initialization, an assignment, an argument, a returned value -- and the widening records which member it is, as the tag the union object carries.

Going back is not a conversion. The member held is only known at runtime, so recovering it is a test, and a test is what a pattern is: naming a member binds it with its own type when the tag says the union holds it, whether the pattern is written as an arm of a \tokennolst{match} or as an \tokennolst{if let}. Nothing about a union value says statically which member it is, so no conversion can produce one.

The base a union declares is no exception. \tokennolst{union U over B} says every member descends from \tokennolst{B}, which is what lets the methods of \tokennolst{B} be called on the union directly; it does not make the union an instance of \tokennolst{B}. The union object is its own aggregate -- a tag, a vtable and the overlay of the members -- and the \tokennolst{\&B} is inside that overlay, reached the same way as any other member.

\subsubsection*{Example}

\textit{test\_\allowbreak{}resources/\allowbreak{}type\_\allowbreak{}union/\allowbreak{}test37.yr}:

%% check: skip
\begin{lstlisting}[style=coloredverbatimError]
class A {
    pub self () {}
}

class B {
    pub self () {}
}

union U
| A
| B;

fn take (u : &U)-> &A {
    u
}
\end{lstlisting}

\begin{lstlisting}[style=bashVerb, breaklines=true, escapechar=~]
Error[E4313] : union &(test37::U) cannot be converted to &(test37::A)
 --> test_resources/type_union/test37.yr:(15,20)
15  ~\lstbox{\textSFxi{}}~ fn take (u : &U)-> &A {
    ~\lstbox{\textSFv{}}~                    ^
16  ~\lstbox{\textSFxi{}}~     u
    ~\lstbox{\textSFv{}}~     -
 --> test_resources/type_union/test37.yr:(11,7)
11  ~\lstbox{\textSFxi{}}~ union U
    ~\lstbox{\textSFv{}}~       -  note: the member held is only known at runtime
    ~\lstbox{\textSFxi{}}~
    = when validating test37::take -- (test_resources/type_union/test37.yr:(15,4))
    ~\lstbox{\textSFii{}}~~\lstbox{\textSFx{}}~~\lstbox{\textSFx{}}~~\lstbox{\textSFx{}}~~\lstbox{\textSFx{}}~~\lstbox{\textSFx{}}~~\lstbox{\textSFx{}}~~\lstbox{\textSFx{}}~
\end{lstlisting}

\subsubsection*{How to fix it}

Name the member in a pattern -- an arm of a \tokennolst{match} over the union, or an \tokennolst{if let} -- and work there, which is where the member has its own type.

\errorcode{E4314}{union \errhole{} does not list the member \errhole{}}
\label{sec:codes:e4314}

Raised during \textbf{validation}, from \texttt{ValidateErrorMessage::\allowbreak{}UNION\_\allowbreak{}MEMBER\_\allowbreak{}UNKNOWN} (\textit{src/\allowbreak{}ymirc/\allowbreak{}semantic/\allowbreak{}validator/\allowbreak{}errors.yr}).

\subsubsection*{What it means}

A pattern over a union names the member it tests, and the union holds one of the members it declares and nothing else. A type the union does not list can never be the one held, so the arm naming it matches nothing and is rejected rather than left as dead code.

The member is matched by identity, not by convertibility: a type that merely converts to a declared member is not that member either.

\subsubsection*{Example}

\textit{test\_\allowbreak{}resources/\allowbreak{}type\_\allowbreak{}union/\allowbreak{}test48.yr}:

%% check: skip
\begin{lstlisting}[style=coloredverbatimError]
record X {
    pub self () {}
}

record Y {
    pub self () {}
}

record Z {
    pub self () {}
}

union U
| X
| Y;

fn main () {
    let u : U = X ();
    match u {
        z : Z => { z; }
        _ => {}
    }
}
\end{lstlisting}

\begin{lstlisting}[style=bashVerb, breaklines=true, escapechar=~]
Error[E4314] : union test48::U does not list the member test48::Z
 --> test_resources/type_union/test48.yr:(22,13)
22  ~\lstbox{\textSFxi{}}~         z : Z => { z; }
    ~\lstbox{\textSFv{}}~             ^
 --> test_resources/type_union/test48.yr:(15,7)
15  ~\lstbox{\textSFxi{}}~ union U
    ~\lstbox{\textSFv{}}~       -
    ~\lstbox{\textSFxi{}}~
    = when validating pattern matching with value u of type test48::U -- (test_resources/type_union/test48.yr:(21,5))
    = when validating test48::main -- (test_resources/type_union/test48.yr:(19,4))
    ~\lstbox{\textSFii{}}~~\lstbox{\textSFx{}}~~\lstbox{\textSFx{}}~~\lstbox{\textSFx{}}~~\lstbox{\textSFx{}}~~\lstbox{\textSFx{}}~~\lstbox{\textSFx{}}~~\lstbox{\textSFx{}}~
\end{lstlisting}

\subsubsection*{How to fix it}

Name a member the union declares, or add the type to the union when it is meant to be one.

\errorcode{E4315}{\errhole{} is not a union type}
\label{sec:codes:e4315}

Raised during \textbf{validation}, from \texttt{ValidateErrorMessage::\allowbreak{}NOT\_\allowbreak{}A\_\allowbreak{}UNION} (\textit{src/\allowbreak{}ymirc/\allowbreak{}semantic/\allowbreak{}validator/\allowbreak{}errors.yr}).

\subsubsection*{What it means}

A template parameter written \tokennolst{union T} constrains the kind of type that may be bound to it, the way \tokennolst{class T} or \tokennolst{tuple T} do: it accepts a union and nothing else. The constraint is checked wherever the argument comes from -- written explicitly at the call site, or deduced from a parameter typed \tokennolst{T}.

What fails is the candidate, not the program: a template whose parameter cannot be bound is discarded, so this diagnostic is reported as the reason under the error naming the call, and the call fails only when no candidate is left.

\subsubsection*{Example}

\textit{test\_\allowbreak{}resources/\allowbreak{}type\_\allowbreak{}union/\allowbreak{}test62.yr}:

%% check: skip
\begin{lstlisting}[style=coloredverbatimError]
fn keep {union T} (u : T)-> T {
    u
}

fn main () {
    keep (12);
}
\end{lstlisting}

\begin{lstlisting}[style=bashVerb, breaklines=true, escapechar=~]
Error[E4226] : the call operator is not defined for test62::keep {union T}(u: T)-> T and {12: i32}
 --> test_resources/type_union/test62.yr:(8,10)
 8  ~\lstbox{\textSFxi{}}~     keep (12);
    ~\lstbox{\textSFv{}}~          ^
 --> test_resources/type_union/test62.yr:(3,4)
 3  ~\lstbox{\textSFxi{}}~ fn keep {union T} (u : T)-> T {
    ~\lstbox{\textSFv{}}~    ----  note: candidate test62::keep {union T}(u: T)-> T
    ~\lstbox{\textSFxi{}}~       ~\lstbox{\textSFii{}}~~\lstbox{\textSFx{}}~ error: i32 is not a union type
    ~\lstbox{\textSFxi{}}~
    = when validating test62::main -- (test_resources/type_union/test62.yr:(7,4))
    ~\lstbox{\textSFii{}}~~\lstbox{\textSFx{}}~~\lstbox{\textSFx{}}~~\lstbox{\textSFx{}}~~\lstbox{\textSFx{}}~~\lstbox{\textSFx{}}~~\lstbox{\textSFx{}}~~\lstbox{\textSFx{}}~
\end{lstlisting}

\subsubsection*{How to fix it}

Pass a union, or drop the \tokennolst{union} constraint when the parameter is meant to accept any type.

\errorcode{E4317}{union \errhole{} does not implement trait \errhole{}}
\label{sec:codes:e4317}

Raised during \textbf{validation}, from \texttt{ValidateErrorMessage::\allowbreak{}UNION\_\allowbreak{}NOT\_\allowbreak{}IMPL} (\textit{src/\allowbreak{}ymirc/\allowbreak{}semantic/\allowbreak{}validator/\allowbreak{}errors.yr}).

\subsubsection*{What it means}

A union is used where the trait named is required, and the union does not implement it.

A union implements a trait when every one of its members implements it, so that the methods of the trait can be dispatched to whichever member the union holds. A union declaring a base with \tokennolst{over} is dispatched through the vtable of that base, so it implements only the traits the base implements, even when every member implements more. The notes name the members, or the base, that do not implement the trait.

\subsubsection*{Example}

\textit{test\_\allowbreak{}resources/\allowbreak{}type\_\allowbreak{}union/\allowbreak{}test69.yr}:

%% check: skip
\begin{lstlisting}[style=coloredverbatimError]
// A union having a member that does not implement a trait does not implement it.

use std::stream;

record A {
    impl Streamable;
}

record C {}

union Y
| A
| C;

fn foo {T impl Streamable} (_ : T) {}

fn bar (y : Y) {
    foo (y);
}
\end{lstlisting}

\begin{lstlisting}[style=bashVerb, breaklines=true, escapechar=~]
Error[E4226] : the call operator is not defined for test69::foo {T impl Streamable}(_: T) and {y: test69::Y}
 --> test_resources/type_union/test69.yr:(18,9)
18  ~\lstbox{\textSFxi{}}~     foo (y);
    ~\lstbox{\textSFv{}}~         ^
 --> test_resources/type_union/test69.yr:(15,4)
15  ~\lstbox{\textSFxi{}}~ fn foo {T impl Streamable} (_ : T) {}
    ~\lstbox{\textSFv{}}~    ---  note: candidate test69::foo {T impl Streamable}(_: T)
    ~\lstbox{\textSFxi{}}~      ~\lstbox{\textSFviii{}}~~\lstbox{\textSFx{}}~ error: union test69::Y does not implement trait Streamable
    ~\lstbox{\textSFxi{}}~      ~\lstbox{\textSFii{}}~~\lstbox{\textSFx{}}~ note: the member test69::C does not implement trait Streamable
    ~\lstbox{\textSFxi{}}~
    = when validating test69::bar -- (test_resources/type_union/test69.yr:(17,4))
    ~\lstbox{\textSFii{}}~~\lstbox{\textSFx{}}~~\lstbox{\textSFx{}}~~\lstbox{\textSFx{}}~~\lstbox{\textSFx{}}~~\lstbox{\textSFx{}}~~\lstbox{\textSFx{}}~~\lstbox{\textSFx{}}~
\end{lstlisting}

\subsubsection*{How to fix it}

Implement the trait in every member of the union -- or, for a union declaring a base, in the base itself.
